\newpage
\recipe{Apple Pie}{1 h 15 min}{8 servings}{Classic family favorite}

\begin{multicols}{2}
	\subsection*{Ingredients}
	\begin{ingredients}
		\item 1 double pie crust\footnote{You can use store-bought shortcuts or a homemade all-butter crust.}
		\item 6 apples, peeled and sliced
		\item 150 g sugar
		\item 1 tsp cinnamon
		\item 1 tbsp lemon juice\footnote{Or 2 if, you like lemon.}
		\item 2 tbsp butter, cubed
		\item 1 egg, beaten
	\end{ingredients}
	
	\columnbreak
	
	\subsection*{You'll also need}
	\begin{equipment}
		\item 9-inch pie dish
		\item Large mixing bowl
		\item Apple peeler \& corer
		\item Pastry brush
		\item Rolling pin
	\end{equipment}
\end{multicols}

\medskip

% --- Preparation ---
Preheat the oven to 190\,\textdegree C (375\,\textdegree F). In a large mixing bowl, gently toss the peeled and sliced apples with the sugar, cinnamon, and lemon juice until they are evenly coated. Line your pie dish with the first rolled-out crust and carefully pile the apple mixture inside, ensuring it's evenly distributed. 

Dot the top of the fruit with the cubed\footnote{Or just sliced} butter, then lay the second crust over the top. Crimp the edges tightly to seal in the juices, cut a few small slits in the center to let steam escape, and brush the entire top crust with the beaten egg wash. Bake for 45 to 50 minutes until the pastry is deeply golden brown, then let it cool completely on a wire rack before slicing.

\vspace*{0.75 cm}

% --- Description ---
\begin{recipeblock}
	\noindent\textbf{The Story behind this Pie:} This recipe has been passed down through generations. The secret is letting the apples macerate just long enough to draw out the juices without making the bottom crust soggy.
\end{recipeblock}